\documentclass[11pt,a4paper]{article}
\usepackage[utf8]{inputenc}
\usepackage[T1]{fontenc}
\usepackage[margin=2.5cm]{geometry}
\usepackage{amsmath,amssymb,bm}
\usepackage{booktabs}
\usepackage{array}
\usepackage[dvipsnames]{xcolor}
\usepackage[colorlinks=true,linkcolor=Navy,urlcolor=Navy]{hyperref}
\definecolor{Navy}{RGB}{27,73,101}
\usepackage{mdframed}
\usepackage{titlesec}
\titleformat{\section}{\normalfont\large\bfseries\color{Navy}}{\thesection}{0.7em}{}
\setlength{\emergencystretch}{2.5em}

\newcommand{\ket}[1]{\lvert #1 \rangle}
\newcommand{\bra}[1]{\langle #1 \rvert}
\newcommand{\dg}{^{\dagger}}
\newcommand{\Om}{\Omega}
\newcommand{\Omu}{\Omega_{\mu}}
\newcommand{\Dmu}{\Delta_{\mu}}

\newmdenv[linecolor=Navy,linewidth=0.8pt,backgroundcolor=Navy!4,
          roundcorner=3pt,innertopmargin=7pt,innerbottommargin=7pt]{keybox}

\title{\bfseries From four levels to the Dicke model\\[3pt]
\large a step-by-step reduction for cavity Rydberg-EIT electrometry}
\author{}
\date{\today}

\begin{document}
\maketitle
\vspace{-10mm}

\begin{abstract}\noindent
We start from the exact laboratory Hamiltonian and reduce it, one controlled step at a time,
to a collective spin--boson model. Each step is a matrix operation and each carries a stated
validity condition. Two results deserve advance notice. First, the enormous
$\ket{r}\!\to\!\ket{r'}$ splitting does not appear in the reduced model at all: the rotating
frame absorbs it, leaving only the microwave \emph{detuning} $\Dmu$. Second, the microwave
dressing generally produces two branches that couple to the ground state \emph{unequally}, by
an amount set by the mixing angle; the familiar symmetric Autler--Townes doublet is the
$\Dmu=0$ limit, and the single-branch (AC Stark) picture is the $|\Dmu|\gg\Omu$ limit. Which
limit one works in decides whether the measurand enters linearly or quadratically. Throughout
$\hbar=1$.
\end{abstract}

\section{The exact Hamiltonian}

Four atomic levels with laboratory energies $\omega_1=0<\omega_2<\omega_3<\omega_4$:
\[
  \ket{1}=\ket{5S_{1/2}},\quad \ket{2}=\ket{5P_{3/2}},\quad
  \ket{3}=\ket{53D_{5/2}},\quad \ket{4}=\ket{54P_{3/2}},
\]
numerically $\omega_2/2\pi=384.25$, $\omega_3/2\pi=1008.82$~THz, and
$(\omega_4-\omega_3)/2\pi\simeq10$~GHz. Writing $\sigma_{jk}=\ket{j}\bra{k}$, and with no
approximations whatsoever,
\begin{equation}
\begin{split}
  H = {}& \underbrace{\sum_{j=1}^{4}\omega_j\,\sigma_{jj} + \omega_{\rm cav}\,a\dg a}_{\text{free}}
   + \underbrace{g_0\,(a+a\dg)(\sigma_{12}+\sigma_{21})}_{\text{cavity, }\ket{1}\leftrightarrow\ket{2}} \\[2ex]
   & + \underbrace{\Om_c\cos(\nu_c t)\,(\sigma_{23}+\sigma_{32})}_{\text{control laser}}
   + \underbrace{\Omu\cos(\nu_\mu t)\,(\sigma_{34}+\sigma_{43})}_{\text{microwave}} .
\end{split}
  \label{eq:lab}
\end{equation}
The classical fields carry their full time dependence; the cavity coupling carries both
$a$ and $a\dg$. Nothing has been dropped.

\section{Step 1: the rotating frame}

Transform with $\tilde H = U H U\dg + i\dot U U\dg$, choosing
\begin{equation}
  U(t)=\exp\Bigl\{i t\bigl[\omega_L\, a\dg a + \omega_L\,\sigma_{22}
        + (\omega_L+\nu_c)\,\sigma_{33} + (\omega_L+\nu_c+\nu_\mu)\,\sigma_{44}\bigr]\Bigr\},
\end{equation}
i.e.\ each level is rotated at the \emph{sum of the drive frequencies needed to reach it}
($\omega_L$ is a reference near the cavity resonance). This is an exact change of picture. Its
whole effect is on the diagonal:

\begin{center}\small
\begin{tabular}{@{}llll@{}}
\toprule
level & rotated at & lab energy & rotating-frame energy \\
\midrule
$\ket{1}$ & $0$ & $0$ & $0$\\
$\ket{2}$ & $\omega_L$ & $2\pi\!\times\!384$~THz & $\Delta \equiv \omega_2-\omega_L$ \quad($\sim$GHz)\\
$\ket{3}$ & $\omega_L+\nu_c$ & $2\pi\!\times\!1009$~THz & $\delta \equiv \omega_3-\omega_L-\nu_c$ \quad($\sim$MHz)\\
$\ket{4}$ & $\omega_L+\nu_c+\nu_\mu$ & $\omega_3+2\pi\!\times\!10$~GHz &
   $\delta+\Dmu$, \ \ $\Dmu \equiv (\omega_4-\omega_3)-\nu_\mu$\\
\bottomrule
\end{tabular}
\end{center}

\begin{keybox}
\textbf{This answers the obvious objection.} $\ket{3}$ and $\ket{4}$ are \emph{not} degenerate
--- they are $10$~GHz apart. But the frame rotates $\ket{4}$ faster than $\ket{3}$ by exactly
$\nu_\mu$, and $\nu_\mu$ is chosen close to that splitting. The $10$~GHz is subtracted off;
what remains on the diagonal is only the residual \emph{microwave detuning} $\Dmu$, of order
MHz. Setting $\Dmu=0$ makes the two entries equal --- and that is a statement about the
\emph{drive}, not about the atom.
\end{keybox}

\section{Step 2: the rotating-wave approximation}

In the new frame the cavity term $a\dg\sigma_{21}$ picks up $e^{2i\omega_L t}$, and
$\cos(\nu t)\sigma_{jk}$ splits into a static piece and one at $e^{2i\nu t}$. Dropping the
fast pieces requires
\begin{equation}
  g_0,\ \Om_c,\ \Omu \ \ll\ \omega_L,\ \nu_c,\ \nu_\mu ,
  \tag{A1}
\end{equation}
which is spectacularly well satisfied (MHz against THz, and MHz against the $10$~GHz microwave).
The result is time-independent and still exact to this order:
\begin{equation}
  \tilde H = \omega_c\, a\dg a +
  \begin{pmatrix}
    0        & g_0 a\dg & 0        & 0 \\[2pt]
    g_0 a    & \Delta   & \Om_c/2  & 0 \\[2pt]
    0        & \Om_c/2  & \delta   & \Omu/2 \\[2pt]
    0        & 0        & \Omu/2   & \delta+\Dmu
  \end{pmatrix},
  \qquad \omega_c \equiv \omega_{\rm cav}-\omega_L .
  \label{eq:H4}
\end{equation}
Note the matrix is tridiagonal: each field connects only neighbouring levels.

\section{Step 3: eliminate the intermediate state}

Detune far from the lossy state $\ket{2}$,
\begin{equation}
  |\Delta| \ \gg\ \Om_c,\ g_0\sqrt{\langle a\dg a\rangle},\ \Gamma_e ,
  \tag{A2}
\end{equation}
so its amplitude follows adiabatically: $c_2 \simeq -(g_0 a\,c_1 + \tfrac{\Om_c}{2}c_3)/\Delta$.
Substituting into Eq.~\eqref{eq:H4} leaves, in $\{\ket{1},\ket{3},\ket{4}\}$,
\begin{equation}
  \begin{pmatrix}
    -\dfrac{g_0^2 a\dg a}{\Delta} & -\lambda\,a\dg & 0 \\[9pt]
    -\lambda\,a & \delta - \dfrac{\Om_c^2}{4\Delta} & \Omu/2 \\[9pt]
    0 & \Omu/2 & \delta+\Dmu
  \end{pmatrix},
  \qquad
  \boxed{\ \lambda = \frac{g_0\Om_c}{2\Delta}\ }
  \label{eq:H3}
\end{equation}
The off-diagonal $\lambda a$ is the two-photon Raman coupling we want. The diagonal entries are
light shifts: $-\Om_c^2/4\Delta$ on $\ket{3}$, absorbed by redefining $\delta$, and a
photon-number-dependent shift on $\ket{1}$, negligible for
$g_0^2\langle a\dg a\rangle\ll|\Delta|\,\omega_c$. The overall sign of $\lambda$ is a phase
convention and we drop it.

\begin{keybox}
\textbf{Where the dominant systematic comes from.} $\lambda\propto\Om_c$ and the light shift
$\propto\Om_c^2$ both track the control-laser intensity, and the shift adds \emph{directly} to
$\delta$ --- the same channel the measurand uses. Control-laser intensity noise is therefore
structurally indistinguishable from signal.
\end{keybox}

\section{Step 4: dress the Rydberg pair}

The lower-right block of Eq.~\eqref{eq:H3} is a two-level problem on its own:
\begin{equation}
  H_\mu=\begin{pmatrix}\delta & \Omu/2\\ \Omu/2 & \delta+\Dmu\end{pmatrix}
      = \Bigl(\delta+\tfrac{\Dmu}{2}\Bigr)\mathbf{1}
        + \tfrac12\bigl(\Omu\,\sigma_x + \Dmu\,\sigma_z\bigr).
\end{equation}
Diagonalise with the rotation $R(\theta)$ defined by
\begin{equation}
  \tan\theta = \frac{\Omu}{\Dmu},
  \qquad
  \ket{+} = \sin\tfrac\theta2\ket{3}+\cos\tfrac\theta2\ket{4},
  \qquad
  \ket{-} = \cos\tfrac\theta2\ket{3}-\sin\tfrac\theta2\ket{4},
\end{equation}
giving the \emph{generalised Rabi} splitting
\begin{equation}
  \boxed{\ \omega_\pm=\delta+\frac{\Dmu}{2}\pm\frac{W}{2},
  \qquad W=\sqrt{\Omu^{2}+\Dmu^{2}}\ }
  \label{eq:branches}
\end{equation}
Now rotate the rest of Eq.~\eqref{eq:H3}. The Raman coupling connects $\ket{1}$ to the
\emph{bare} $\ket{3}=\sin\tfrac\theta2\ket{+}+\cos\tfrac\theta2\ket{-}$, so it splits between
the branches in that ratio. In the basis $(\ket{g},\ket{+},\ket{-})$:
\begin{equation}
  \boxed{\;
  H^{(3)}=\begin{pmatrix}
    0 & g_+ a\dg & g_- a\dg\\[3pt]
    g_+ a & \omega_+ & 0\\[3pt]
    g_- a & 0 & \omega_-
  \end{pmatrix},
  \qquad
  g_+=\lambda\sin\tfrac\theta2,\quad g_-=\lambda\cos\tfrac\theta2 \;}
  \label{eq:core}
\end{equation}
There is no $\ket{+}\!\leftrightarrow\!\ket{-}$ entry: that transition is the microwave one and
carries no cavity matrix element.

\subsection*{The two limits, and why they are the whole story}

\begin{center}\small
\begin{tabular}{@{}>{\raggedright\arraybackslash}p{3.5cm}>{\raggedright\arraybackslash}p{5.0cm}>{\raggedright\arraybackslash}p{5.6cm}@{}}
\toprule
& \textbf{Resonant:} $\Dmu=0$ & \textbf{Far detuned:} $|\Dmu|\gg\Omu$\\
\midrule
mixing angle & $\theta=\pi/2$ & $\theta\to0$\\
branches & $\ket{\pm}=(\ket{3}\pm\ket{4})/\sqrt2$ & $\ket{-}\to\ket{3}$, $\ket{+}\to\ket{4}$\\
splitting $W$ & $\Omu$ \ (linear in $E_\mu$) & $|\Dmu|+\Omu^2/2|\Dmu|$\\
couplings & $g_\pm=\lambda/\sqrt2$ \ (equal) & $g_-\to\lambda$, $g_+\to0$ \ (one branch)\\
shift of the coupled level & $\pm\Omu/2$, \textbf{linear} & $-\Omu^2/4\Dmu$, \textbf{quadratic}\\
this is & Autler--Townes & AC Stark\\
used by & this proposal & Ding \emph{et al.}, Wang \emph{et al.}\\
\bottomrule
\end{tabular}
\end{center}

\noindent One mixing angle therefore separates the two ways of doing Rydberg electrometry. On
resonance the measurand enters as a level \emph{splitting}, linear in $E_\mu$ with the exact
dipole scale factor, at the price of two branches that both couple. Far detuned it enters as a
\emph{shift}, quadratic in $E_\mu$ with a polarizability-dependent factor, but the model is
genuinely two-level. There is no setting that gives both.

\section{Step 5: many atoms, and the Dicke form}

Take $N$ atoms in $\ket{g}$ and work at low excitation, $n_++n_-\ll N$ (A3). Holstein--Primakoff
replaces each collective excitation by a boson $b_\pm$, and the couplings add constructively,
$\lambda\to\lambda\sqrt N$. Write $G_\pm = \sqrt N\,g_\pm$. Equation~\eqref{eq:core} then gives
\begin{equation}
  H=\omega_c a\dg a+\sum_{s=\pm}\Bigl[\omega_s b_s\dg b_s + G_s\bigl(a\dg b_s+a b_s\dg\bigr)\Bigr],
\end{equation}
which conserves excitation number: starting from $\ket{g,0}$ nothing ever happens, since it is
an exact eigenstate. A Dicke model requires the counter-rotating terms $a\dg b_s\dg + ab_s$ as
well, supplied by a second drive tone (or one modulated tone). With both,
\begin{equation}
  \boxed{\;H=\omega_c a\dg a+\sum_{s=\pm}\Bigl[\omega_s b_s\dg b_s
      + G_s\,(a+a\dg)(b_s+b_s\dg)\Bigr]\;}
  \label{eq:dicke}
\end{equation}

\section{Step 6: the threshold}

In quadratures $x_j=(o_j+o_j\dg)/\sqrt2$ the potential energy is $\tfrac12\bm{x}^{\!\top}V\bm{x}$,
\begin{equation}
  V=\begin{pmatrix}
      \omega_c & 2G_+ & 2G_-\\
      2G_+ & \omega_+ & 0\\
      2G_- & 0 & \omega_-
    \end{pmatrix}.
\end{equation}
The normal phase (all atoms in $\ket{g}$, cavity dark) is stable while $V$ is positive definite,
and loses stability at $\det V=0$:
\begin{equation}
  \boxed{\ \omega_c\,\omega_+\omega_- \;=\; 4\bigl(G_+^2\,\omega_- + G_-^2\,\omega_+\bigr)\ }
  \label{eq:threshold}
\end{equation}
Positive definiteness also demands $\omega_->0$; if the lower branch falls below $\ket{g}$ the
normal phase fails for a mundane reason --- a two-photon resonance --- rather than by symmetry
breaking. With $G=\lambda\sqrt N$ the two limits read
\begin{align}
  \Dmu=0:&\qquad \omega_c\,\omega_+\omega_- = 2G^2(\omega_++\omega_-)=4G^2\delta,\\
  |\Dmu|\gg\Omu:&\qquad \omega_c\,\omega_- = 4G^2 ,
\end{align}
the second being the textbook two-level Dicke threshold $G^{\rm c}=\tfrac12\sqrt{\omega_c\omega_-}$
for a spin of splitting $\omega_-$. \emph{This} is the effective two-level model --- and it is
legitimate only in the far-detuned limit, where the measurand has become quadratic.

\section{Summary}

\begin{center}\small
\begin{tabular}{@{}>{\raggedright\arraybackslash}p{3.4cm}>{\raggedright\arraybackslash}p{3.6cm}>{\raggedright\arraybackslash}p{7.0cm}@{}}
\toprule
\textbf{Step} & \textbf{Condition} & \textbf{Effect}\\
\midrule
Rotating frame & exact & subtracts the drive frequencies; the $10$~GHz
  $\ket{3}$--$\ket{4}$ splitting becomes $\Dmu$\\
RWA & $g_0,\Om_c,\Omu\ll\omega_L,\nu_c,\nu_\mu$ & removes all time dependence; Eq.~\eqref{eq:H4}\\
Eliminate $\ket{2}$ & $|\Delta|\gg\Om_c,g_0\sqrt{\bar n},\Gamma_e$ &
  $\lambda=g_0\Om_c/2\Delta$, plus light shifts (the dominant systematic)\\
Dress $\ket{3},\ket{4}$ & exact & $\omega_\pm=\delta+\Dmu/2\pm W/2$, couplings
  $\lambda\sin\tfrac\theta2$, $\lambda\cos\tfrac\theta2$\\
Holstein--Primakoff & $n_++n_-\ll N$ & $\lambda\to G=\lambda\sqrt N$\\
Counter-rotating & second drive tone & without it $\ket{g,0}$ is an eigenstate and nothing radiates\\
Two-level limit & $|\Dmu|\gg\Omu$ & one branch decouples; measurand becomes quadratic\\
\bottomrule
\end{tabular}
\end{center}

\noindent Only the last step is optional, and it is the one that decides what kind of sensor
this is.

\end{document}
